% Source: source/普通高考/2018/2018全国2理(甘肃,青海,内蒙古,黑龙江,吉林,辽宁,海南.宁夏,新疆,陕西,重庆).pdf, page 1.
% The source flowchart is vector line art; pdfimages -list found no embedded bitmap.
% Elements: start/end rounded rectangles, assignment rectangles, decision diamond,
% arithmetic process rectangles, output parallelogram, directed connectors, and branch labels.
% Relations: 开始 -> N=0,T=0 -> i=1 -> i<100; 是 -> N update -> T update -> blank increment
% -> loop back to the decision; 否 -> S=N-T -> 输出 S -> 结束.
% solid segments: all box outlines and connectors; dashed segments: none.
% labels: 是 at the left decision branch and 否 at the right branch; math labels centered
% in boxes, with the left return path outside the process column to avoid overlap.
% Node-edge list from the PDF: start→init, init→seti, seti→check; check(是)→addn,
% check(否)→sub; addn→addt, addt→increment, increment→check; sub→out, out→finish.
\begin{tikzpicture}[
  >=Stealth,
  line width=0.45pt,
  line join=round,
  line cap=round,
  font=\normalsize,
  flowtext/.style={anchor=center, align=center,
                   execute at begin node={%
                     \everypar{}%
                     \setlength{\parindent}{0pt}%
                     \setlength{\leftskip}{0pt}%
                     \setlength{\rightskip}{0pt}%
                     \setlength{\parfillskip}{0pt}%
                   }},
  every node/.style={flowtext},
  flowbox/.style={draw, flowtext, minimum height=0.72cm, inner xsep=5pt, inner ysep=3pt},
  process/.style={flowbox, minimum width=3.05cm},
  smallprocess/.style={flowbox, minimum width=1.55cm},
  startstop/.style={flowbox, rounded corners=0.18cm, minimum width=1.60cm},
  output/.style={flowbox, trapezium, trapezium left angle=75, trapezium right angle=105,
                minimum width=3.00cm},
  decision/.style={flowbox, diamond, aspect=2.8, minimum width=3.80cm, minimum height=1.18cm},
]
  % Preserve the source's clear white margin around the outer return path.
  \path[use as bounding box] (-5.15,-0.12) rectangle (4.05,8.82);

  % Coordinates preserve the schematic topology and the source's two-column flow.
  \node[startstop] (start) at (0,8.40) {开始};
  \node[process] (init) at (0,7.20) {$N=0, T=0$};
  \node[smallprocess] (seti) at (0,5.98) {$i=1$};
  \node[decision] (check) at (0,4.65) {$i<100$};

  \node[process] (addn) at (-2.35,3.25) {$N=N+\dfrac{1}{i}$};
  \node[process, minimum width=3.55cm] (addt) at (-2.35,1.78) {$T=T+\dfrac{1}{i+1}$};
  \node[smallprocess, minimum width=2.45cm] (increment) at (-2.35,0.35) {};

  \node[process, minimum width=3.30cm] (sub) at (2.25,3.25) {$S=N-T$};
  \node[output] (out) at (2.25,1.86) {输出 $S$};
  \node[startstop] (finish) at (2.25,0.55) {结束};

  % Main vertical path.
  \draw[->] (start.south) -- (init.north);
  \draw[->] (init.south) -- (seti.north);
  % The loop return joins the central path just above the decision.
  \coordinate (loopjoin) at (0,5.38);
  \draw (seti.south) -- (loopjoin);

  % Branches from the decision diamond.
  \draw[->] (check.west) -- (-2.35,4.65) -- (addn.north)
    node[pos=0.20, above=1pt] {是};
  \draw[->] (check.east) -- (2.25,4.65) -- (sub.north)
    node[pos=0.20, above=1pt] {否};

  % Left loop body and return path.
  \draw[->] (addn.south) -- (addt.north);
  \draw[->] (addt.south) -- (increment.north);
  \draw[->] (increment.west) -- (-4.85,0.35) -- (-4.85,5.38) -- (loopjoin);
  \draw[->] (loopjoin) -- (check.north);

  % False branch output path.
  \draw[->] (sub.south) -- (out.north);
  \draw[->] (out.south) -- (finish.north);
\end{tikzpicture}
